\documentclass[10pt,a4paper]{amsart}
\usepackage{amsmath,amssymb,amsthm,graphicx}
% Retained minimal preamble for the Stacks Project chapter more-algebra.tex (More on Algebra).
% The chapter file pinned-original-source.tex carries only \input{preamble} on its first line, so no
% preamble text lives inside the chapter itself.  The Stacks Project uses one single shared preamble for
% every chapter; the snapshot stages it identically as sources/stacks-chow/preamble.tex and
% sources/stacks-etale-cohomology/preamble.tex.  The block below is copied mechanically, line for line,
% from sources/stacks-chow/preamble.tex lines 14-18 (the Xy-pic setup) and 154-217 (the native theorem
% declarations together with the mathematical macros).  Nothing was added, removed or redefined; the
% wrapper class and the amsmath/amssymb/amsthm/graphicx base packages are supplied by the assembler.
% Chapter: more-algebra.tex -- More on Algebra.
% For commutative diagrams we use Xy-pic
\usepackage[all]{xy}

% We use 2cell for 2-commutative diagrams.
\xyoption{2cell}

% Theorem environments.
%
\theoremstyle{plain}
\newtheorem{theorem}[subsection]{Theorem}
\newtheorem{proposition}[subsection]{Proposition}
\newtheorem{lemma}[subsection]{Lemma}

\theoremstyle{definition}
\newtheorem{definition}[subsection]{Definition}
\newtheorem{example}[subsection]{Example}
\newtheorem{exercise}[subsection]{Exercise}
\newtheorem{situation}[subsection]{Situation}

\theoremstyle{remark}
\newtheorem{remark}[subsection]{Remark}
\newtheorem{remarks}[subsection]{Remarks}

\numberwithin{equation}{subsection}

% Macros
%
\def\lim{\mathop{\mathrm{lim}}\nolimits}
\def\colim{\mathop{\mathrm{colim}}\nolimits}
\def\Spec{\mathop{\mathrm{Spec}}}
\def\Hom{\mathop{\mathrm{Hom}}\nolimits}
\def\Ext{\mathop{\mathrm{Ext}}\nolimits}
\def\SheafHom{\mathop{\mathcal{H}\!\mathit{om}}\nolimits}
\def\SheafExt{\mathop{\mathcal{E}\!\mathit{xt}}\nolimits}
\def\Sch{\mathit{Sch}}
\def\Mor{\mathop{\mathrm{Mor}}\nolimits}
\def\Ob{\mathop{\mathrm{Ob}}\nolimits}
\def\Sh{\mathop{\mathit{Sh}}\nolimits}
\def\NL{\mathop{N\!L}\nolimits}
\def\CH{\mathop{\mathrm{CH}}\nolimits}
\def\proetale{{pro\text{-}\acute{e}tale}}
\def\etale{{\acute{e}tale}}
\def\QCoh{\mathit{QCoh}}
\def\Ker{\mathop{\mathrm{Ker}}}
\def\Im{\mathop{\mathrm{Im}}}
\def\Coker{\mathop{\mathrm{Coker}}}
\def\Coim{\mathop{\mathrm{Coim}}}

% Boxtimes
%
\DeclareMathSymbol{\boxtimes}{\mathbin}{AMSa}{"02}

%
% Macros for moduli stacks/spaces
%
\def\QCohstack{\mathcal{QC}\!\mathit{oh}}
\def\Cohstack{\mathcal{C}\!\mathit{oh}}
\def\Spacesstack{\mathcal{S}\!\mathit{paces}}
\def\Quotfunctor{\mathrm{Quot}}
\def\Hilbfunctor{\mathrm{Hilb}}
\def\Curvesstack{\mathcal{C}\!\mathit{urves}}
\def\Polarizedstack{\mathcal{P}\!\mathit{olarized}}
\def\Complexesstack{\mathcal{C}\!\mathit{omplexes}}
% \Pic is the operator that assigns to X its picard group, usage \Pic(X)
% \Picardstack_{X/B} denotes the Picard stack of X over B
% \Picardfunctor_{X/B} denotes the Picard functor of X over B
\def\Pic{\mathop{\mathrm{Pic}}\nolimits}
\def\Picardstack{\mathcal{P}\!\mathit{ic}}
\def\Picardfunctor{\mathrm{Pic}}
\def\Deformationcategory{\mathcal{D}\!\mathit{ef}}
\usepackage{cleveref}
\crefname{originallemma}{original lemma}{original lemmas}
\Crefname{originallemma}{Original lemma}{Original lemmas}
\makeatletter
\newlabel{algebra-lemma-finite-flat-local}{{00NZ}{source}{}{sourcepointer-1}{}}
\newlabel{algebra-lemma-finite-flat-local@cref}{{[originallemma][0][]00NZ}{[1][0][]source}{}{}{}}
\makeatother
\crefname{originallemma}{original lemma}{original lemmas}
\Crefname{originallemma}{Original lemma}{Original lemmas}
\makeatletter
\newlabel{algebra-lemma-finite-projective}{{00NX}{source}{}{sourcepointer-2}{}}
\newlabel{algebra-lemma-finite-projective@cref}{{[originallemma][0][]00NX}{[1][0][]source}{}{}{}}
\makeatother
\crefname{originallemma}{original lemma}{original lemmas}
\Crefname{originallemma}{Original lemma}{Original lemmas}
\makeatletter
\newlabel{algebra-lemma-NAK}{{00DV}{source}{}{sourcepointer-3}{}}
\newlabel{algebra-lemma-NAK@cref}{{[originallemma][0][]00DV}{[1][0][]source}{}{}{}}
\makeatother
\crefname{originallemma}{original lemma}{original lemmas}
\Crefname{originallemma}{Original lemma}{Original lemmas}
\makeatletter
\newlabel{algebra-lemma-fun}{{05G8}{source}{}{sourcepointer-4}{}}
\newlabel{algebra-lemma-fun@cref}{{[originallemma][0][]05G8}{[1][0][]source}{}{}{}}
\makeatother
\title{Isomorphism classes of finite projective modules over a ring are detected modulo an ideal in the Jacobson radical}
\author{The Stacks Project Authors}
\date{stacks-project:more-algebra.tex:local-snapshot-20261012; GNU FDL 1.2 or later}
\begin{document}
\maketitle
\setcounter{subsection}{1}
\noindent\textit{Editorial note.} All delivered mathematical text is the author's own, copied from the pinned Stacks Project chapter file more-algebra.tex; no statement or proof step was rewritten, reordered or removed.\par\medskip
\noindent\textit{Editorial note.} The chapter file pinned-original-source.tex carries only \textbackslash{}input\{preamble\} on its first line, so no preamble text lives inside the chapter. The retained file sources/stacks-more-algebra/retained-preamble.tex is a mechanical, line-for-line copy of the project's single shared preamble as staged in the snapshot, namely sources/stacks-chow/preamble.tex lines 14-18 (Xy-pic setup) and 154-217 (theorem declarations and mathematical macros). Nothing was added, removed or redefined, no package was added, and the excerpt serves only as preamble, never as a span coordinate.\par\medskip
\noindent\textit{Editorial note.} The source URL is the source-level Stacks Project URL, as required when more than one card is produced from one source; the per-claim tags are carried by claim\_key and by the external references. The same-chapter pointer lemma-lift-stably-free (0BC5) is counted by a sibling card of this batch; the sentence that refers to it sits at source line 141 and is deliberately kept outside both counted spans.\par\medskip
\noindent\textit{Allowed standard prerequisites (editorial).}\begin{itemize}
\item Nakayama's lemma for a finitely generated module over a ring with an ideal contained in the Jacobson radical. Algebra, Lemma algebra-lemma-NAK (canonical Stacks tag 00DV), invoked at proof lines 162-163. Turns the surjectivity of the two composites modulo I into genuine surjectivity of psi o phi and phi o psi as maps of R-modules.
\item A surjective endomorphism of a finite module is an isomorphism. Algebra, Lemma algebra-lemma-fun (canonical Stacks tag 05G8), invoked at proof line 164. Upgrades the surjectivity obtained from Nakayama to invertibility of the two composites, which is what makes phi an isomorphism.
\item The lifting property of projective modules against a surjection of rings, applied to finite projective modules. the projectivity of P and P' recorded in the statement at source lines 148-149, used at proof lines 159-161 and 168-170. Produces the lift psi : P' -> P of the composite P' -> P'/IP' -> P/IP and the lift phi : P -> P' of P -> P/IP -> P'/IP'; both are the objects the remainder of the proof operates on.
\end{itemize}\medskip
\ifdefined\hypertarget\hypertarget{sourcepointer-1}{}\fi
\noindent\textit{Source reference.} Original lemma 00NZ: focused-seeds/AGP/sources/stacks-snapshot-20261011/tags.txt, source lines 877--877. Cross-chapter pointer into the Algebra chapter; it is the standard source of a reduction and the argument built on top of that reduction stays inside the counted proof span.\par\medskip
\ifdefined\hypertarget\hypertarget{sourcepointer-2}{}\fi
\noindent\textit{Source reference.} Original lemma 00NX: focused-seeds/AGP/sources/stacks-snapshot-20261011/tags.txt, source lines 875--875. Cross-chapter pointer into the Algebra chapter supplying a standard criterion of finite projectivity; it is background terminology, not part of the counted argument.\par\medskip
\ifdefined\hypertarget\hypertarget{sourcepointer-3}{}\fi
\noindent\textit{Source reference.} Original lemma 00DV: focused-seeds/AGP/sources/stacks-snapshot-20261011/tags.txt, source lines 512--512. Cross-chapter pointer into the Algebra chapter for Nakayama's lemma, a standard prerequisite; both applications of it and the surrounding argument stay inside the counted proof span.\par\medskip
\ifdefined\hypertarget\hypertarget{sourcepointer-4}{}\fi
\noindent\textit{Source reference.} Original lemma 05G8: focused-seeds/AGP/sources/stacks-snapshot-20261011/tags.txt, source lines 7052--7052. Cross-chapter pointer into the Algebra chapter for the fact that a surjective endomorphism of a finite module is an isomorphism, used only to pass from surjectivity to invertibility.\par\medskip
\begin{lemma}
\label{lemma-isomorphic-finite-projective-lifts}
Let $R$ be a ring. Let $I \subset R$ be an ideal. Assume that
every element of $1 + I$ is a unit (in other words $I$ is contained
in the Jacobson radical of $R$). If $P$ and $P'$ are finite
projective $R$-modules, then
\begin{enumerate}
\item if $\varphi : P \to P'$ is an $R$-module map inducing an
isomorphism $\overline{\varphi} : P/IP \to P'/IP'$, then $\varphi$
is an isomorphism,
\item if $P/IP \cong P'/IP'$, then $P \cong P'$.
\end{enumerate}
\end{lemma}

\begin{proof}
Proof of (1). As $P'$ is projective as an $R$-module we may
choose a lift $\psi : P' \to P$ of the map
$P' \to P'/IP' \xrightarrow{\overline{\varphi}^{-1}} P/IP$.
By Nakayama's lemma (Algebra, Lemma \ref{algebra-lemma-NAK})
$\psi \circ \varphi$ and $\varphi \circ \psi$ are surjective.
Hence these maps are isomorphisms (Algebra, Lemma \ref{algebra-lemma-fun}).
Thus $\varphi$ is an isomorphism.

\medskip\noindent
Proof of (2). Choose an isomorphism $P/IP \cong P'/IP'$.
Since $P$ is projective we can choose a lift $\varphi : P \to P'$ of the map
$P \to P/IP \to P'/IP'$. Then $\varphi$ is an isomorphism by (1).
\end{proof}

\end{document}
